%% More Strategies, Same Zero — IEEE conference format (IEEEtran).
%% Compile on Overleaf: New Project -> Upload -> this file -> Recompile (menu: pdfLaTeX).
%% Or locally: pdflatex false_alpha_ieee.tex (x2).  Class IEEEtran ships with Overleaf/TeXLive.
\documentclass[conference]{IEEEtran}
\IEEEoverridecommandlockouts
\usepackage{amsmath,amssymb,amsfonts}
\usepackage{booktabs}
\usepackage{array}
\usepackage{url}
\usepackage[hidelinks]{hyperref}
\def\BibTeX{{\rm B\kern-.05em{\sc i\kern-.025em b}\kern-.08em
    T\kern-.1667em\lower.7ex\hbox{E}\kern-.125emX}}

\begin{document}

\title{More Strategies, Same Zero: Does LLM-Scale Alpha Search\\
Discover Edge or Manufacture False Discoveries?}

\author{\IEEEauthorblockN{Parv Mehndiratta}
\IEEEauthorblockA{Independent Researcher\\
whizkid783@gmail.com}}

\maketitle

\begin{abstract}
Large language models (LLMs) can propose and backtest trading strategies at effectively
unbounded scale, reviving backtest overfitting in a regime where the number of trials $N$ is
enormous and cheap. We ask whether scaling generative strategy search accumulates genuine
edge or only false discoveries that vanish under proper multiple-testing correction. Using a
dependency-free alpha engine with a strict verifier (causal next-bar execution, transaction
costs, worst-of-regimes out-of-sample Sharpe, and a deflated-Sharpe haircut), we sweep $N$
from 10 to 3{,}000 across three generators---uniform-random, evolutionary ``creative''
synthesis, and a live LLM (Claude Opus~4.8)---on four controls (strong/small planted edge,
pure noise, real data). The naive best-of-$N$ Sharpe rises with $N$ on every market including
pure noise; the gate holds survivors at exactly \textbf{zero} on the real S\&P~500 at all $N$
and all generators, while recovering a strong planted edge. A \emph{power surface} shows the
detectable-edge threshold rising $\sim\!\sqrt{2\ln N}$, so moderate real edges become
unprovable at scale. The zero generalizes under a distribution-free block bootstrap, on a
real cross-sectional equity panel, and across FX, gold, and crypto---seven real markets, zero
survivors. We derive the \textbf{worst-of-regimes deflation} the gate actually needs (the
standard single-Sharpe deflated Sharpe over-deflates the min-over-regimes statistic by
$\sim\!2.4\times$) and a \textbf{robustness-dividend identity}: requiring worst-of-$k$
independent regimes is equivalent to a single-Sharpe search over only $N'\!\approx\!7$
strategies, an implicit multiple-testing correction that collapses as regimes correlate
(measured: $\rho_\mathrm{reg}\!=\!0.67$ on real S\&P, $N'\!\approx\!84$). We report four
refutations, including two analytic effective-trial-count shortcuts that manufacture false
discoveries and are corrected by the bootstrap. All results are reproducible from committed
scripts; a verification suite of 20 adversarial checks and 38 unit tests is green.
\end{abstract}

\begin{IEEEkeywords}
multiple testing, backtest overfitting, deflated Sharpe ratio, large language models,
data snooping, quantitative finance, false discovery, extreme value theory
\end{IEEEkeywords}

\section{Introduction}
A backtest certifies the past, not the future, and the more strategies one tests the more
certain it is that the best is lucky rather than skilled. Bailey and L\'opez de Prado~\cite{dsr}
formalize this as the false-strategy theorem: under the null of no skill, the expected maximum
Sharpe ratio over $N$ independent trials grows like $\sqrt{2\ln N}$, so any sufficiently large
search ``discovers'' an impressive-looking strategy. The classical remedy is to \emph{deflate}
the observed Sharpe by exactly this expected best-of-$N$ (the Deflated Sharpe Ratio, DSR).

Large language models change the operative regime of this theorem: they propose candidate
strategies essentially without limit and with superficial novelty, so $N$---historically
bounded by human and compute budgets---becomes large, cheap, and correlated. Recent systems
explicitly mine alphas with LLMs~\cite{alphaagent} and report in-sample or single-holdout
gains. This raises a direct, testable question:
\emph{when a generative engine proposes and backtests more and more trading strategies, do
genuinely profitable ones accumulate, or only statistically false ones that vanish under
proper multiple-testing correction?}

\textbf{Contributions.} The statistic is not new (the DSR is~\cite{dsr}, the Reality Check is
White's~\cite{rc}, and the ``LLM edge evaporates'' conclusion is established by
FINSABER~\cite{finsaber}). Our contributions are: (1) a \emph{derived worst-of-regimes
deflation}---the gate scores the minimum over $k$ regimes, yet the DSR is the expected max of
$N$ \emph{single} Sharpes; we derive the correct envelope $\mathbb{E}[\max_N \min_k Z]$ and
show the DSR over-deflates it by ${\sim}2.4\times$ (Sec.~\ref{sec:deflation}); (2) a
\emph{robustness-dividend identity} linking regime robustness to multiple testing, which
collapses as regimes correlate; (3) a controlled, multi-generator (including a live LLM),
seven-market, power-surface protocol with planted/noise calibration; and (4) four reported
negatives, including two analytic effective-trial-count shortcuts that manufacture false
discoveries and are corrected by the bootstrap.

\section{The Gate and the N-Sweep}\label{sec:method}
A strategy (``alpha'') is an expression in a small S-expression DSL over causal price/volume
features, combined by unary/binary operators (depth $\le 5$); it evaluates causally to a
per-bar signal squashed to a position in $[-1,1]$. The \textbf{verifier} (the gate) is
adversarial: (i) no look-ahead, $\mathrm{strat}[t]=\mathrm{pos}[t{-}1]\,r[t]-c\,|\Delta
\mathrm{pos}|$; (ii) transaction cost $c=10$\,bps per unit turnover; (iii) the timeline is
split into an in-sample segment and $k$ contiguous out-of-sample (OOS) regimes, and a strategy
is scored by its \emph{worst} regime; (iv) the score is the worst-regime Sharpe minus the
expected best-of-$N$ under the null (the DSR haircut). A strategy is ``verified'' only if its
annualized, after-cost, worst-regime, deflated OOS Sharpe clears a positive bar
($0.5$).

Three generators stand in for the search engine: \texttt{random} (uniform draws over the DSL,
$\approx$ i.i.d.\ brute force), \texttt{creative} (Boden-style operators over a shared
substrate~\cite{boden}), and a \texttt{live LLM} (Claude Opus~4.8). For each (market,
generator) we draw a pool, backtest each alpha once, and apply the deflation analytically per
$N$. We report, versus $N$: the best naive pooled-OOS Sharpe (the data-snoop curve);
\texttt{robust}, the count clearing worst-regime$+$cost but \emph{not} deflation (the
ablation); and \texttt{survivors}, the count clearing the full deflated gate. Four markets
form the matrix: a strong planted edge (positive control), a realistically small planted edge,
a pure random-walk null, and real data.

\section{Results: The Zero, and Why It Is Not Trivial}
\textbf{The naive metric is fooled, worse as $N$ grows.} The expected best-of-$N$ pooled-OOS
Sharpe rises monotonically on \emph{every} market, including pure noise ($+0.04\!\to\!+0.34$)
and real S\&P~500 ($+0.72\!\to\!+0.99$ annualized); a naive practitioner would accept the
${\sim}1.0$ Sharpe ``winner.''

\textbf{The deflated gate holds at zero on real data, and recovers a real edge.} On the real
S\&P~500 the survivor count is \textbf{0 at every $N$}, for all three generators (Table~\ref
{tab:surv}). This is not a degenerate ``reject everything'': on the strong-planted market the
\emph{same} gate passes strategies, and passes more as $N$ grows (Table~\ref{tab:surv}),
because the edge is genuinely there.

\begin{table}[t]\centering
\caption{Gate survivors vs.\ $N$ (creative generator).}
\label{tab:surv}
\begin{tabular}{lcccc}
\toprule
market & $N{=}10$ & $100$ & $1000$ & $3000$\\
\midrule
real S\&P 500 & 0.0 & 0.0 & 0.0 & \textbf{0.0}\\
noise (null)  & 0.0 & 0.0 & 0.0 & 0.0\\
planted (edge)& 0.7 & 7.1 & 58.3 & \textbf{154.4}\\
\bottomrule
\end{tabular}
\end{table}

\textbf{Ablation: which tightener does the work.} On the adversarial noise market the
worst-regime requirement \emph{alone} kills the snoop (\texttt{robust}$=0$). On real data it
does not: because 2016--2024 sub-periods resemble each other, \texttt{robust} grows with $N$
to 115--196 ``regime-robust'' mirages; the \emph{multiple-testing deflation} is what collapses
these to zero. On real data the deflation is load-bearing.

\textbf{The live-LLM arm.} Claude Opus~4.8 reaches $N{=}300$ on real S\&P: best-of-$N$ rises
$+0.48\!\to\!+0.80$, survivors $=0$, mean return correlation $\bar\rho=+0.007$. Of four
pre-registered hypotheses, three hold; the hypothesis that the LLM is \emph{more} correlated
than random is \textbf{refuted} ($\bar\rho$ no higher than random's $+0.029$). A new finding
cuts the other way: the LLM's structural diversity saturates---1{,}000 draws yield only 160
distinct valid expressions (vs.\ ${\sim}280$ for random)---so scaling an LLM proposer buys
less new search than scaling random generation.

\textbf{The power surface (central result).} A gate that returns 0 on real data is only
interesting if it \emph{would} return $>0$ for a realistic edge. Sweeping the planted edge
magnitude $\times$ search scale $N$ (Table~\ref{tab:power}), a genuine edge admitted at
$N{=}10$ becomes undetectable at $N\!\ge\!30$: the same true edge, now rejected, because
searching more raised the bar above it. The detectable-edge threshold rises monotonically with
$N$ ($\approx\sqrt{2\ln N}$, the empirical instantiation of~\cite{dsr}). Consequently we do
\emph{not} claim market efficiency; we claim that \textbf{no edge large enough to survive a
search of this scale exists in this DSL/data}---a high-precision statement with a known,
large Type-II region. Inverting the gate gives the minimum detectable edge: under the
corrected deflation the recall floor rises $0.69\,(N{=}10)\!\to\!1.48\,(1000)\!\to\!1.62\,
(3000)$, while real markets offer a best worst-regime Sharpe of ${\approx}0.85$---below the
floor. We measure the precision side directly by held-out replication: every survivor on one
sample of the strong-planted market also survives an independent sample, so the
false-discovery rate is ${\approx}0$ even at $N{=}1000$; the price is recall, not precision.

\begin{table}[t]\centering
\caption{Power surface: survivors over (edge strength $\times$ $N$).}
\label{tab:power}
\begin{tabular}{lcccccc}
\toprule
edge $\backslash$ $N$ & 10 & 30 & 100 & 300 & 1000 & 3000\\
\midrule
$0.00$--$0.15$ & 0 & 0 & 0 & 0 & 0 & 0\\
$0.25$ & \textbf{1} & 0 & 0 & 0 & 0 & 0\\
$0.40$ (strong) & 2 & 2 & 8 & 20 & 61 & 168\\
\bottomrule
\end{tabular}
\end{table}

\textbf{Generalization.} The zero is not a parametric artifact: a stationary block
bootstrap~\cite{politis} feeding White's Reality Check / Romano--Wolf~\cite{rc,romanowolf} and
Hansen's SPA~\cite{spa}, with the worst-of-regimes operator preserved inside each resample,
gives 0 survivors on real data at every $N$ and block length. The result also holds on a real
15-stock cross-sectional equity panel (range/volume features active) and across FX, gold, and
crypto---seven real markets, zero survivors at every $N$.

\section{The Worst-of-Regimes Deflation and the Robustness Dividend}\label{sec:deflation}
The DSR is the expected maximum of $N$ \emph{single} Sharpes, but our gate scores the
\emph{minimum over $k$ regimes}, then the best of $N$. The correct null is therefore
$\mathbb{E}\big[\max_{i\le N}\min_{g\le k} Z_{i,g}\big]$. Because the minimum of $k$ standard
normals has a thin upper tail, its best-of-$N$ envelope sits near the
$\big(1-N^{-1/k}\big)$ quantile, $\Phi^{-1}\!\big(1-N^{-1/k}\big)$. Monte-Carlo validation
gives an envelope $2.3$--$2.8\times$ smaller than the single-Sharpe one at relevant $N$ (e.g.,
$3.26\!\to\!1.38$ at $N{=}1000,k{=}3$): \textbf{the standard DSR over-deflates the
worst-regime statistic by ${\sim}2.4\times$.} Re-running with the corrected, far weaker bar,
the best deflated Sharpe on real markets rises from ${\approx}{-}1.5$ to ${\approx}{-}0.3$ but
\emph{stays negative}---0 survivors. The zero is robust to the correct deflation, not an
over-deflation artifact.

This yields a unifying identity. Define the deflation-equivalent search size $N'$ such that a
single-Sharpe search over $N'$ has the same null maximum as worst-of-$k$ at $N$. For
\emph{independent} regimes, worst-of-3 at $N{=}1000$ gives $N'\!\approx\!7$: \textbf{requiring
a strategy to survive several independent regimes is an implicit multiple-testing correction}
as strong as shrinking a thousand-strategy search to a handful. Generalizing to equicorrelated
regimes with correlation $\rho_\mathrm{reg}$, $N'$ climbs $7\!\to\!28\!\to\!62\!\to\!160\!\to\!
426$ as $\rho_\mathrm{reg}\!:\,0\!\to\!0.95$. We then \emph{measure} $\rho_\mathrm{reg}$ as the
mean cross-sectional correlation between strategies' per-regime Sharpes: the adversarial
synthetic null sits at $\rho_\mathrm{reg}\!=\!0.24$ ($N'\!\approx\!16$, dividend intact),
while real S\&P~500 sits at $\rho_\mathrm{reg}\!=\!0.67$ ($N'\!\approx\!84$, a ${\sim}12\times$
collapse). This resolves the ablation: on the synthetic null the dividend protects (so
worst-regime alone suffices), while on real data the dividend has collapsed (so the explicit
deflation must do the work)---robustness and multiple testing trade off quantitatively along
regime correlation.

\section{Reported Negatives}
In the project's ``report what does not work'' discipline, we record four refutations.
(i)~The pre-registered hypothesis that the LLM searches more redundantly than random is
refuted in return space ($\bar\rho$). (ii)--(iii)~Two analytic effective-trial-count
shortcuts fail: importing the genomics effective-number-of-independent-tests
machinery~\cite{nyholt,liji} via the participation ratio gives $M_\mathrm{eff}\!\approx\!5$,
a near-zero haircut, and ${\sim}500$ false ``survivors'' on real markets that the bootstrap
rejects; an envelope-inversion estimate is autocorrelation-confounded. Pursuing the cause with
a serial-correlation-adjusted Sharpe standard error~\cite{lo} shows the Newey--West inflation
is negligible ($\eta^2\!\approx\!1.07$), yet the cross-sectional variance of the standardized
Sharpe is ${\approx}4.6$, not 1: the real cause is \textbf{structural heterogeneity}
(${\sim}78\%$ of the cross-sectional Sharpe variance is between-strategy), so a generative
search produces $N$ structurally different strategies, not $N$ i.i.d.\ draws of one null---the
``$N$ i.i.d.\ trials'' premise fails on the \emph{identical} part. The block bootstrap, which
resamples each strategy's own return path, is therefore the necessary arbiter, not a mere
check. (iv)~A pre-registered prediction that the false-discovery rate rises with $N$ is
refuted (it stays ${\approx}0$): family-wise control here yields ${\approx}$zero false
discoveries for free.

\section{Related Work}
The closest direct precedent is FINSABER~\cite{finsaber}, which finds LLM trading agents
generate no statistically significant alpha under broad, long-horizon evaluation and concludes
LLM-derived alpha is largely a methodological artefact; it controls snooping by evaluation
design, not by a formal multiple-testing gate, an $N$-sweep, a power surface, or planted
controls. The statistical machinery is classical: White's Reality Check~\cite{rc}, the
stationary bootstrap~\cite{politis}, Sullivan--Timmermann--White's data-snooping correction
over a 7{,}846-rule universe (which already reported a zero on the S\&P)~\cite{stw}, Hansen's
SPA~\cite{spa}, Romano--Wolf step-down~\cite{romanowolf}, and the deflated Sharpe~\cite{dsr}.
The factor-zoo literature~\cite{harvey} studies multiple testing in cross-sectional asset
pricing. Same-lineage LLM alpha-mining~\cite{alphaagent} targets discovery/decay, not
multiple-testing correction. Our worst-of-regimes deflation and the effective-number-of-tests
framing draw on the genomics correlated-testing literature~\cite{nyholt,liji}, where linkage
disequilibrium between SNPs is the analogue of shared structure between strategies; we report
(Sec.~V) that the naive import fails for the maximum statistic.

\section{Limitations and Conclusion}
The headline caveat is low recall: ``zero survivors'' means ``no edge provable at this search
scale,'' never ``no edge exists.'' The DSL and data are modest (a small daily feature set);
counts are finite-pool bootstrap means; the live LLM reaches $N{=}300$ and a single model.
Within those bounds, the conclusion is robust across seven real markets, a distribution-free
test, eight random seeds, and the corrected deflation. \emph{LLM-scale strategy search
manufactures false discoveries, and the only correction that stops them simultaneously raises
the bar for proving any genuine edge---so beyond a point, additional search yields only
mirages and unprovable maybes.} Robustness across regimes and multiple-testing correction are
two faces of the same coin; how much each contributes depends, quantitatively, on regime
correlation.

\section*{Reproducibility}
Every reported number is produced by a committed script and written to a JSON artifact; runs
are CRC-seeded and deterministic. Code, data-acquisition scripts, a 20-check adversarial
verification suite, and 38 unit tests are available in the Mentat repository.

\begin{thebibliography}{99}
\bibitem{dsr} D.~H. Bailey and M.~L\'opez de Prado, ``The deflated Sharpe ratio: correcting
for selection bias, backtest overfitting, and non-normality,'' \emph{J. Portfolio
Management}, vol.~40, no.~5, pp.~94--107, 2014.
\bibitem{rc} H.~White, ``A reality check for data snooping,'' \emph{Econometrica}, vol.~68,
no.~5, pp.~1097--1126, 2000.
\bibitem{stw} R.~Sullivan, A.~Timmermann, and H.~White, ``Data-snooping, technical trading
rule performance, and the bootstrap,'' \emph{J. Finance}, vol.~54, no.~5, pp.~1647--1691, 1999.
\bibitem{spa} P.~R. Hansen, ``A test for superior predictive ability,'' \emph{J. Business \&
Economic Statistics}, vol.~23, no.~4, pp.~365--380, 2005.
\bibitem{romanowolf} J.~P. Romano and M.~Wolf, ``Stepwise multiple testing as formalized data
snooping,'' \emph{Econometrica}, vol.~73, no.~4, pp.~1237--1282, 2005.
\bibitem{harvey} C.~R. Harvey, Y.~Liu, and H.~Zhu, ``$\ldots$ and the cross-section of expected
returns,'' \emph{Rev. Financial Studies}, vol.~29, no.~1, pp.~5--68, 2016.
\bibitem{politis} D.~N. Politis and J.~P. Romano, ``The stationary bootstrap,'' \emph{J.
American Statistical Association}, vol.~89, no.~428, pp.~1303--1313, 1994.
\bibitem{lo} A.~W. Lo, ``The statistics of Sharpe ratios,'' \emph{Financial Analysts Journal},
vol.~58, no.~4, pp.~36--52, 2002.
\bibitem{nyholt} D.~R. Nyholt, ``A simple correction for multiple testing for single-nucleotide
polymorphisms in linkage disequilibrium with each other,'' \emph{American J. Human Genetics},
vol.~74, no.~4, pp.~765--769, 2004.
\bibitem{liji} J.~Li and L.~Ji, ``Adjusting multiple testing in multilocus analyses using the
eigenvalues of a correlation matrix,'' \emph{Heredity}, vol.~95, no.~3, pp.~221--227, 2005.
\bibitem{finsaber} % VERIFY exact authors/title before submission
  Z.~Liang \emph{et al.}, ``FINSABER: rigorous backtesting reveals the fragility of
  LLM-based trading strategies,'' arXiv:2505.07078, 2025.
\bibitem{alphaagent} % VERIFY author list before submission
  Z.~Tang \emph{et al.}, ``AlphaAgent: LLM-driven alpha mining with regularized exploration to
  counteract alpha decay,'' arXiv:2502.16789, 2025 (KDD~2025).
\bibitem{boden} M.~A. Boden, \emph{The Creative Mind: Myths and Mechanisms}, 2nd~ed. London,
U.K.: Routledge, 2004.
\bibitem{afml} M.~L\'opez de Prado, \emph{Advances in Financial Machine Learning}. Hoboken,
NJ, USA: Wiley, 2018.
\end{thebibliography}

\end{document}
